% ===== 原文 PDF pp.83--85，并补完 p.86 前 3 行（印刷 pp.82--84）=====
% [边界] 起点：承接 tr07b 文件末的半句「…考虑式 (2.50)，Wess--Zumino 贡献」。p.83 页首即该句的
%   后半，故本文件开头直接续写，不重译上文。
% [原文 p.82]
$R(A)$ 与 $R(B)$ 由奇点形变理论确定如下 [48]。这里 $R(A)$ 由 CB 算子的标度维数决定，即由 $\Delta[\lambda_{\alpha}]>1$ 的 $\lambda_{\alpha}$ 决定；而 $R(B)$ 从 $\hat F$ 的余维一奇点处获得贡献——在这些奇点处出现一个额外的无质量超多重态，且局域贡献由垂直于该奇点的局域坐标的 $U(1)_r$ 荷给出。对于由 IIB 型弦论中 IHS 实现的 4d SCFT\footnote{在这类 4d SCFT 中也没有自由超多重态，即 (2.50) 中 $h=0$。}，它们由 [48] 显式给出为
\begin{equation}
  R(A) = \sum_{\Delta[\lambda_{\alpha}]>1}\Delta[\lambda_{\alpha}]-r\,,\qquad
  R(B) = \frac{\hat\alpha_{max}}{4}\,\mu\,,
  \tag{5.18}
\end{equation}
其中 $\hat\alpha_{max}$ 指 CB 谱中最大的标度维数，它由
\begin{equation}
  \hat\alpha_{max} = \left(\sum_{i}^{4}q_i-1\right)^{-1}
  \tag{5.19}\end{equation}
给出。

% [Tutor Note: (5.18)(5.19) 要放回 (2.50) 才有意义：$a=\frac{5r}{24}+\frac{h}{24}+\frac{R(A)}{4}+\frac{R(B)}{6}$，
%   $c=\frac{r}{6}+\frac{h}{12}+\frac{R(B)}{3}$；本类理论 $h=0$（脚注 59）。注意 (5.19) 的 $\hat\alpha_{max}$
%   就是 (5.11) 的 $a$：二者都等于 $1/(\sum_i q_i-1)$，这并非巧合——CB 上最大的标度维数正好是
%   「荷为 1 的形变参数」的维数，而 (5.11) 正是用 $\Delta(\Omega_3)=1$ 定出的比例常数。]
% [Tutor Note: $R(B)$ 里出现的是 $\mu$ 而 $R(A)$ 里出现的是 CB 谱：因为 $R(A)$ 只由 $\Delta>1$ 的 VEV 参数
%   （即库仑分支）贡献，而 $R(B)$ 统计的是形变三重复形 $\hat F$ 上余维一奇点处额外无质量超多重态
%   的贡献。$\mu$ 恰是这种奇点的计数（每个单形变参数 $\lambda_\alpha$ 对应一处），所以 $R(B)\propto\mu$。]

\SUB{5.2}{从 $(G,G')$ 型奇点得到 $\cN=2$ SCFT 的例子}

本节中，让我们聚焦于一类特殊的孤立超曲面奇点，称为 $(G,G')$ 型奇点。在 [72] 中，作者们考虑了这样一类孤立三重复形超曲面奇点：其定义多项式由两个定义 ADE 曲面奇点的多项式之和给出，
\begin{equation}
  F:\quad F_{G}(x_1,x_2)+F_{G'}(x_3,x_4) = 0\,,
  \tag{5.20}
\end{equation}
其中每个 $F_{G}$ 属于如下类型之一：
\begin{equation}
  \begin{aligned}
    &F_{A_k}(x,y) = x^2+y^{k+1}\,,\\
    &F_{D_k}(x,y) = x^2y+y^{k-1}\,,\\
    &F_{E_6}(x,y) = x^3+y^{4}\,,\\
    &F_{E_7}(x,y) = x^3+xy^{3}\,,\\
    &F_{E_8}(x,y) = x^3+y^{5}\,.
  \end{aligned}
  \tag{5.21}
\end{equation}

% [Tutor Note: (5.21) 是 ADE 曲面奇点的标准（拟齐次）正规形，注意它与 (5.9) 的写法并不相同：
%   (5.9) 里的 $F_{A_k}$ 是 $\C^4$ 中三重复形超曲面（四个变量），而 (5.21) 的 $F_G$ 只是 $\C^2$ 中曲面
%   奇点的正规形（两个变量）。两者通过 (5.20) 的求和相联系：$G$ 与 $G'$ 各贡献一对变量。
%   因此 Milnor 数因子化：$\mu=\mu_G\,\mu_{G'}$；而对 ADE 曲面奇点 $\mu_G=r_G$（例如 $A_n$ 的
%   $\mu=n=r_{A_n}$），于是 $\mu=r_Gr_{G'}$，这正是脚注 60 中 $\mu=2r+f=r_Gr_{G'}$ 的来源。]

事实证明，这类奇点会产生一大类 Argyres--Douglas 理论，被称为 $(G,G')$ 理论\footnote{注意，这类 AD 理论之所以标注为 $(G,G')$，是因为与之对应的 BPS 箭图具有 $G$ 与 $G'$ 的 Dynkin 图之积的形状。特别地，$\mu=2r+f=r_Gr_{G'}$。}。注意标号 $(G,G')$ 并不唯一：由 (5.20) 很容易看出存在显然的同构 $(G,G')\sim(G',G)$。此外，还
% [原文 p.83]
存在许多进一步的等价关系 [47,72]，例如
\begin{equation}
  (A_1,D_4)\sim(A_2,A_2)\,,\quad (A_1,E_6)\sim(A_3,A_2)\,,\quad (D_4,A_3)\sim(E_6,A_2)\,,\quad (E_8,A_3)\sim(E_6,A_4)\,,\ \dots
  \tag{5.22}
\end{equation}

% [Tutor Note: (5.22) 的这些等价关系都来自坐标的显式重定义：把 $F_G+F_{G'}$ 的某几项配方/换元后，
%   会得到另一对 $(G_1,G_2)$ 的求和形式；由于 $(G,G')$ 只是描述奇点的一种「记号」，奇点本身相同
%   并不意味着标号相同。实用上这意味着：同一个 Argyres--Douglas 理论可以挂在不同的
%   $(G,G')$ 标签下，读文献时不要把它们当成不同的理论。]

\begin{examplebox}{例 5.3：}
  为说明起见，让我们考虑 $(A_1,A_2)$ 奇点的一个显式例子。根据 (5.21)，相应的超曲面由如下多项式定义：
  \begin{equation}
    F = x_1^2+x_2^2+x_3^2+x_4^3\,.
    \tag{5.23}
  \end{equation}
  于是可以很容易读出 $\C^*$ 荷 $q_i=\{\tfrac12,\tfrac12,\tfrac12,\tfrac13\}$ 以及比例常数 $a=\frac{1}{\sum_{i=1}^{4}q_i-1}=\frac{6}{5}$。由 $dF$ 生成的理想是 $(x_1,x_2,x_3,x_4^2)$，因此相应的 Jacobi 代数 $J(F)$ 具有如下单项式基：
  \begin{equation}
    x_4^{\alpha_4}\,,\qquad 0\leqslant\alpha_4\leqslant 1\,.
    \tag{5.24}
  \end{equation}
  然后可以很容易写出相应的 SW 几何，即 (5.6) 的形变
  \begin{equation}
    \hat F = x_1^2+x_2^2+x_3^2+x_4^3+\lambda_1+\lambda_2x_4\,,
    \tag{5.25}
  \end{equation}
  其中 $\lambda_1$ 是标度维数 $\Delta(\lambda_1)=\frac65$ 的 CB VEV，而 $\lambda_2$ 是标度维数 $\Delta(\lambda_2)=\frac45$ 的手征耦合，这与 (5.16) 一致。这正是最初那个由 $\cN=2$ SU(3) 纯 SYM 的红外极限给出的 Argyres--Douglas 理论。

  % [Tutor Note: 用 (5.16) 核一遍这个例子：$\sum_i q_i=\tfrac12+\tfrac12+\tfrac12+\tfrac13=\tfrac{11}{6}$，
  %   故 $a=(11/6-1)^{-1}=6/5$；对 $g=x_4$（$\alpha=(0,0,0,1)$），
  %   $\Delta(\lambda_2)=(1-\tfrac13)/(\tfrac{11}{6}-1)=\tfrac{4}{5}$；对 $g=1$（$\alpha_i=0$），
  %   $\Delta(\lambda_1)=\tfrac{6}{5}$。两者之和 $\tfrac65+\tfrac45=2$，正是 (5.17) 的配对；
  %   而 $\mu=2=2r+f$ 给出 $r=1,\ f=0$——这正是 $(A_1,A_2)$ 理论只有一条库仑分支、
  %   没有味对称性的原因。]
\end{examplebox}

\begin{exercisebox}{习题 5.1}
  计算由上述 $(A_1,A_2)$ 奇点实现的 4d $\cN=2$ SCFT 的中心荷 $(a,c)$。对于一般的 $(A_n,A_m)$ 奇点情形又如何？

  \textbf{提示：}(2.50) 中的 $R(B)$ 在 $(G,G')$ 理论这一大类里约化为 $R(B)=\frac{r_Gr_{G'}}{4}\frac{h^{\vee}_Gh^{\vee}_{G'}}{h^{\vee}_G+h^{\vee}_{G'}}$，其中 $h^{\vee}_G$ 与 $r_G$ 分别是群 $G$ 的对偶 Coxeter 数与秩（参见例如 [211,229]）。
\end{exercisebox}

% [Tutor Note: 习题 5.1 的答案可用来校核本节的公式：$(A_1,A_2)$ 有 $r=1$、$f=0$、$\mu=2$、
%   $\hat\alpha_{max}=\tfrac65$，故 $R(A)=\tfrac65-1=\tfrac15$，$R(B)=\tfrac{1}{4}\cdot\tfrac65\cdot2=\tfrac35$。
%   提示中的公式给出同一结果：$r_{A_1}=1,r_{A_2}=2,h^{\vee}_{A_1}=2,h^{\vee}_{A_2}=3$，
%   于是 $R(B)=\frac{1\cdot2}{4}\cdot\frac{2\cdot3}{2+3}=\frac35$。
%   代回 (2.50)（$h=0$）：$a=\frac{5}{24}+\frac{1}{20}+\frac{1}{10}=\frac{43}{120}$，
%   $c=\frac{1}{6}+\frac{1}{5}=\frac{11}{30}$——正是 $(A_1,A_2)$ Argyres--Douglas 理论的已知值。]

\SUB{5.3}{与 Class S 构造的关系}

到目前为止，我们已经介绍了从孤立超曲面奇点对 4d $\cN=2$ SCFT 作几何工程的一般面貌。人们自然会问：IIB 型工程与第 4 节的 Class S 构造之间是否存在任何关系？答案是有，而且这
% [原文 p.84]
被期待为带五维膜与非平凡内部几何的 M 理论/IIB 对偶的一个推论。

当非紧 Calabi--Yau 流形 $W$ 是 ALE 空间在一条曲线上的纤维化时，我们可以说得更具体。在这种情况下，借助 T-对偶可以显式看到：所得 4d $\cN=2$ SCFT 的库仑分支等价地由一张缠绕黎曼面的五维膜描述，而该黎曼面恰与该 SCFT 的 SW 曲线重合。这一 T-对偶在 [276] 中研究过，它把 ALE 空间 $A_{n-1}$ 奇点附近的 IIB 型理论映射为 IIA 型理论中 $n$ 个重合 NS5 膜组成的堆，而后者的世界体积在低能极限下由 6d $\cN=(2,0)$ 的 $A_{n-1}$ 型 SCFT 描述；同时，IIB 中由缠绕 D3 膜给出的 6d 自对偶弦，对应于终止在 IIA 型 NS5 膜上的 D2 膜的边界。关于该 T-对偶的更多细节，我们请读者参阅综述 [277]。

% [Tutor Note: 这一段是「几何工程 = Class S」的核心机制：$A_{n-1}$ ALE 奇点的 IIB 解与 IIA 中 $n$ 个
%   重合 NS5 膜是 T-对偶的。$A_{n-1}$ 奇点处的 $n$ 个消失 2-圈在 IIA 侧变成 $n$ 张 NS5 膜；
%   NS5 膜的世界体积理论在低能下是 6d $(2,0)$ 的 $A_{n-1}$ 型 SCFT，而要得到 4d 理论就要把它放在
%   一条曲线（Gaiotto 曲线）上紧化。于是「五维膜缠绕黎曼面」与「Class S 的 UV 曲线」是同一件事
%   的两种说法。]

例如，让我们考虑 IIB 型弦论紧化在一个带 $(A_{m-1},A_{n-1})$ 奇点的 Calabi--Yau 流形上：
\begin{equation}
  x^m+y^n+z^2+w^2 = 0\,,
  \tag{5.26}
\end{equation}
它可以进一步看作 ALE 空间的 $A_{m-1}$ 奇点在 $y$-平面上的局部纤维化：
\begin{equation}
  \prod_{i=1}^{m}\bigl(x-a_i(y)\bigr)+z^2+w^2 = 0\,,
  \tag{5.27}
\end{equation}
其中 $a_i(y)$ 可以看作 $y$-平面上的一个多值函数\footnote{与成对的 $(a_i(y),a_j(y))$ 相关联的 ALE 纤维中的非平凡 2-圈，对应于 $A_{m-1}$ 的权。}。由 T-对偶 [276]，这一 IIB 型设置映射为 IIA 型中 $m$ 个 NS5 膜在 $y$-平面上的纤维化堆。把它提升到 M 理论，我们就在 $x=a_i(y)$、$z=w=0$ 处得到 $m$ 个 M5 膜。注意，在 $y=0$ 处，$A_{n-1}$ 奇点出现，$m$ 个 M5 膜会合在一起，它们承载 $\cN=(2,0)$ 的 $A_{m-1}$ 型 SCFT。而在一般的点 $y\neq0$ 处，这 $m$ 个 M5 膜分开，对应于库仑分支上的一个点，4d 有效理论可以由一张缠绕黎曼面 $\Sigma:=\prod_{i=1}^{m}\bigl(x-a_i(y)\bigr)=0$ [26] 的单个 M5 膜的低能极限来描述；该黎曼面可以看作 $y$-平面的一个 $m$ 重覆盖。这就重现了第 4.3 节的 Class S 描述：这里的 Gaiotto 曲线 $\cC$ 被等同于球面 $\mathbb{P}^1$（由 $y$ 局部参数化），并带有相应的非正则穿刺点，如第 4.8.2 与 4.8.3 节所述。

% [Tutor Note: (5.26) 与 (5.21) 的对照：把 $x^m+y^n+z^2+w^2$ 拆成 $(x^m+z^2)+(y^n+w^2)$，
%   两项分别是 $A_{m-1}$（即 $F_{A_k}=x^2+y^{k+1}$ 取 $k+1=m$）与 $A_{n-1}$ 的正规形，
%   故它确实是 $(A_{m-1},A_{n-1})$ 型奇点。这里四个变量的配对方式与 (5.20) 的
%   $F_G(x_1,x_2)+F_{G'}(x_3,x_4)$ 完全相同，只是把 $(x,z)$ 与 $(y,w)$ 各配成一对。
%   读者可用这个拆法把所有 $(A,A)$ 型例子对上 (5.21) 的列表。]

% [Tutor Note: (5.27) 就是 SW 曲线（=Gaiotto 曲线在其 $m$ 重覆盖下的形式）：
%   $\Sigma:\ \prod_{i=1}^m(x-a_i(y))=0$ 的 $m$ 张叶对应 $m$ 个 M5 膜；两张叶相撞的点
%   $a_i(y)=a_j(y)$ 恰是分支点，也即 Gaiotto 曲线上的穿刺点所在。原文把 Gaiotto 曲线 $\cC$
%   取作底空间 $\mathbb{P}^1_y$（而不是 $\Sigma$ 本身），非正则穿刺点位于 $y=0$ 与 $y=\infty$，
%   这与 §4.8.2/§4.8.3 的野穿刺点构造一致。]

此外，IIB 型与 Class S 构造之间通过 T-对偶的这些联系可以推广到所有 $(A,G)$ 理论 [47]，其中 $G$ 可以是上述 ADE 型中的任意一个奇点。为看出这一点，可以取三个坐标（比如 $(x,z,w)$）使它们给出如 (5.21) 中的 $G$ 型奇点 $F_G(x,z,w)=0$（从而由 T-对偶映射为 6d 中 $G$ 型的 $(2,0)$ SCFT），而把另一个坐标（比如 $y$）看作定义 $G$ 型奇点如何纤维化在 $y$-平面之上。

% [Tutor Note: $(A,G)$ 推广的读法：把三个坐标 $(x,z,w)$ 留给 $G$ 型奇点 $F_G=0$（它给出 6d 的
%   $G$ 型 $(2,0)$ 理论），剩下的那一个坐标 $y$ 只负责「纤维化」——于是 4d 理论就是把 6d 的
%   $G$ 型 $(2,0)$ SCFT 放在 $y$-球面 $\mathbb{P}^1$ 上、带一个非正则穿刺点所得的 Class S 理论。
%   取 $G=A_{n-1}$ 就回到 (5.26)/(5.27) 的 $(A_{m-1},A_{n-1})$ 情形；取 $G=D_k,E_k$ 则得到
%   $(A,D)$、$(A,E)$ 型的 Argyres--Douglas 理论。]

% [边界] 本句已写完：p.85 页末「To see that, one can take three coordinates (say (x,z,w)) giving …」
%   跨到 p.86 第 1--3 行「… rise to the type G singularity … fibers over the y-plane.」，按 Lead 裁决
%   由 tr-07 收尾（见 TIMELINE #8）。下一位 tr-08 从 p.86 第 4 行「We thus have a Class S description
%   for the (A,G) theories…」起，不重复本句。
